\documentclass[11pt,a4paper]{article}
\usepackage[utf8]{inputenc}
\usepackage{geometry}
\geometry{
    top=1in,
    bottom=0.5in,
    left=0.5in,
    right=0.5in
}
\usepackage{hyperref}
\usepackage{graphicx}
\usepackage{booktabs}
\usepackage{tabularx}
\usepackage{float}
\usepackage{setspace}
\usepackage{fancyvrb}

\title{\textbf{DSM A2 Part 2 Final Report: MongoDB + Neo4j Querying}}
\author{Shivansh Verma}
\date{\today}

\begin{document}
\maketitle

\tableofcontents
\newpage

\section{MongoDB Querying}

\subsection*{Workflow and Execution}
The full MongoDB Part 2 pipeline is implemented in \texttt{queries/p2\_mongo.py}. It writes complete aggregation pipelines, executes all required queries, saves tabular outputs, and generates visualizations.

Run sequence:
\begin{verbatim}
source .venv/bin/activate
pip install -r requirements.txt
python queries/p2_mongo.py
\end{verbatim}

Generated artifacts:
\begin{itemize}
    \item Pipelines: \texttt{queries/mongo\_part2\_queries.txt}
    \item Data outputs: \texttt{report/p2/data/q*.csv}, \texttt{report/p2/data/q*.json}
    \item Figures: \texttt{report/p2/figures/q*.png}
\end{itemize}

\subsection*{Q1: Cohort Analysis of User Reviewing Behaviour}
Cohorts are defined by year of \texttt{users.yelping\_since}. For each cohort, the query computes mean and standard deviation of stars, mean review character length, mean useful votes per review, and star-bucket proportions.

Key outputs from \texttt{q1\_summary.json}:
\begin{itemize}
    \item Highest mean star cohort: \textbf{2005} (\textbf{3.903})
    \item Highest mean useful-votes cohort: \textbf{2004} (\textbf{2.276})
    \item Total cohorts: \textbf{19}
\end{itemize}

\begin{figure}[H]
    \centering
    \includegraphics[width=0.88\textwidth]{figures/q1_cohort_star_heatmap.png}
    \caption{Q1 heatmap: star proportions by cohort year.}
\end{figure}

\begin{figure}[H]
    \centering
    \includegraphics[width=0.88\textwidth]{figures/q1_cohort_metric_lines.png}
    \caption{Q1 line plot: mean stars and mean useful votes by cohort year.}
\end{figure}

Interpretation:
\begin{itemize}
    \item Older cohorts show stronger useful-vote and rating signals than recent cohorts.
    \item Later cohorts (notably post-2018) trend lower on both dimensions.
\end{itemize}

\subsection*{Q2: Month-over-Month Trend Consistency by Category}
For each category with at least 500 reviews, monthly average stars are computed and consistency is measured as the proportion of consecutive month-pairs with increase/decrease.

Top 3 upward categories:
\begin{itemize}
    \item Japanese Curry (0.544)
    \item Home Cleaning (0.543)
    \item Beer Gardens (0.541)
\end{itemize}

Top 3 downward categories:
\begin{itemize}
    \item Medical Centers (0.566)
    \item Walk-in Clinics (0.554)
    \item Eyelash Service (0.550)
\end{itemize}

\begin{figure}[H]
    \centering
    \includegraphics[width=0.88\textwidth]{figures/q2_trend_consistency_bars.png}
    \caption{Q2 bar plot: consistency ratios for most upward/downward categories.}
\end{figure}

\begin{figure}[H]
    \centering
    \includegraphics[width=0.95\textwidth]{figures/q2_selected_category_monthly_lines.png}
    \caption{Q2 line plot: monthly trajectories for selected categories.}
\end{figure}

Interpretation:
\begin{itemize}
    \item Directional consistency exists, but no category is fully monotonic.
    \item Both upward and downward groups show meaningful month-to-month volatility.
\end{itemize}

\subsection*{Q3: Check-in Class Cross-tabulation}
Businesses are classified into low/medium/high check-in classes using quartiles from check-in counts (Q1=6, Q3=72). For each top-10 category by total review count, metrics are computed by class: mean stars, mean review count, and tips-to-reviews ratio.

\begin{figure}[H]
    \centering
    \includegraphics[width=0.88\textwidth]{figures/q3_tips_reviews_heatmap.png}
    \caption{Q3 heatmap: tips-to-reviews ratio by category and check-in class.}
\end{figure}

\begin{figure}[H]
    \centering
    \includegraphics[width=0.95\textwidth]{figures/q3_mean_star_by_class.png}
    \caption{Q3 bar plot: mean stars by check-in class across top categories.}
\end{figure}

Interpretation:
\begin{itemize}
    \item High check-in businesses generally have substantially higher mean review volumes.
    \item Tips-to-reviews intensity is commonly strongest in the high check-in group.
\end{itemize}

\section{Neo4j / Cypher Querying}

\subsection*{Workflow and Execution}
The full Neo4j Part 2 pipeline is implemented in \texttt{queries/p2\_neo4j.py}. It writes all Cypher/GDS statements to file, executes each section, and stores outputs/figures in report folders.

Run sequence:
\begin{verbatim}
source .venv/bin/activate
pip install -r requirements.txt
python queries/p2_neo4j.py
\end{verbatim}

Generated artifacts:
\begin{itemize}
    \item Cypher/GDS statements: \texttt{queries/cypher\_part2\_queries.txt}
    \item Data outputs: \texttt{report/p2/data/n*.csv}, \texttt{report/p2/data/n*.json}
    \item Figures: \texttt{report/p2/figures/n*.png}
\end{itemize}

\subsection*{NQ1: PageRank on User-Business Projection}
The query projects a directed weighted graph from users to businesses using review stars as weights and runs GDS PageRank with minimum 20 iterations.

Deliverables produced by script:
\begin{itemize}
    \item Top-15 businesses by PageRank with review count and average stars
    \item Spearman correlations: PageRank vs review-count ranking; PageRank vs average-star ranking
    \item Rank-gap analysis for businesses over/under-ranked by PageRank vs simple metrics
\end{itemize}

\begin{figure}[H]
    \centering
    \includegraphics[width=0.92\textwidth]{figures/n1_top15_pagerank_bar.png}
    \caption{NQ1: Top 15 businesses by PageRank score.}
\end{figure}

\begin{figure}[H]
    \centering
    \includegraphics[width=0.78\textwidth]{figures/n1_pagerank_vs_reviewcount_scatter.png}
    \caption{NQ1: PageRank score vs review count.}
\end{figure}

\subsection*{NQ2: Louvain Community Detection}
The script runs Louvain on the \texttt{FRIENDS\_WITH} social graph, filters communities with at least 25 members, and computes top states, top categories, and geographic concentration index (GCI).

\begin{figure}[H]
    \centering
    \includegraphics[width=0.88\textwidth]{figures/n2_top10_geo_concentration.png}
    \caption{NQ2: Top communities by geographic concentration index.}
\end{figure}

\subsection*{NQ3: Node Similarity (Jaccard) for Market Saturation}
Node Similarity is run on business nodes via shared reviewer neighborhoods. Threshold starts at 5 common reviewers and falls back to 4 or 3 if needed. Results include top-5 most saturated and top-5 most fragmented city-category combinations, excluding combinations with fewer than 5 businesses.

\begin{figure}[H]
    \centering
    \includegraphics[width=0.95\textwidth]{figures/n3_market_saturation_top_bottom.png}
    \caption{NQ3: Most saturated and most fragmented city-category combinations.}
\end{figure}

The script also computes mean stars, standard deviation of stars, and mean review count for the most and least saturated combinations.

\subsection*{NQ4: Betweenness vs Degree Centrality}
Betweenness and degree are computed on the social graph. The top-20 overlap is reported, and high-betweenness-low-degree users are compared against high-degree users on review count, distinct cities reviewed, and distinct categories reviewed.

\begin{figure}[H]
    \centering
    \includegraphics[width=0.9\textwidth]{figures/n4_group_comparison_bar.png}
    \caption{NQ4: Group-level comparison for centrality-defined user sets.}
\end{figure}

\subsection*{NQ5: Link Prediction Pipeline for Recommendations}
The script builds a chronological holdout (each user's latest edge as test), runs a GDS link prediction pipeline, and reports AUC-ROC, Precision@10, top predictive features, and top-3 recommendations for 5 sampled users.

\begin{figure}[H]
    \centering
    \includegraphics[width=0.65\textwidth]{figures/n5_lp_metrics_bar.png}
    \caption{NQ5: Link prediction evaluation metrics.}
\end{figure}

\section{Complete Query Statements}

\subsection*{MongoDB Aggregation Pipelines}
\VerbatimInput[fontsize=\scriptsize]{../../queries/mongo_part2_queries.txt}

\subsection*{Neo4j Cypher and GDS Statements}
\VerbatimInput[fontsize=\scriptsize]{../../queries/cypher_part2_queries.txt}

\section{Artifact Index}
\begin{itemize}
    \item Mongo script: \texttt{queries/p2\_mongo.py}
    \item Neo4j script: \texttt{queries/p2\_neo4j.py}
    \item Mongo queries: \texttt{queries/mongo\_part2\_queries.txt}
    \item Neo4j queries: \texttt{queries/cypher\_part2\_queries.txt}
    \item Data outputs: \texttt{report/p2/data/}
    \item Figures: \texttt{report/p2/figures/}
\end{itemize}

\end{document}
